\documentclass[11pt,a4paper]{article}

\usepackage[utf8]{inputenc}

\usepackage[T1]{fontenc}

\usepackage[polish]{babel}

\usepackage{lmodern}

\usepackage{geometry}

\geometry{a4paper,margin=27mm}

\usepackage{microtype}

\usepackage[hidelinks]{hyperref}

% Portfolio: MSC-AR-P01; source: prospective ARC-XR research scope.

\title{Multimodalne obciążenie poznawcze i adaptacja interfejsu AR}

\author{Mateusz Pomianek}

\date{Wersja 0.3.0, 21 listopada 2026}

\begin{document}

\maketitle

\begin{abstract}

Przedmiotem opracowania jest Multimodalne obciążenie poznawcze i adaptacja interfejsu AR. Czy prosta adaptacja zagęszczenia informacji HUD w AR ogranicza błędy względem stałego interfejsu? Przewidziano zrównoważone próby ze stałym HUD oraz z HUD adaptowanym według jednej jawnej reguły; rozbudowany model wielomodalny będzie osobnym rozszerzeniem. Głównym punktem końcowym jest odsetek prób z błędnym wykonaniem zadania, przy jednoznacznym mianowniku. Opis ma charakter wstępny i nie przedstawia niezweryfikowanych wyników.

\end{abstract}

\section{Wprowadzenie}

Rozwój adaptacyjnych systemów XR stwarza możliwość badania interakcji człowieka z informacją przestrzenną, ale samo zastosowanie AR, VR lub modeli sztucznej inteligencji nie przesądza o przewadze poznawczej czy ergonomicznej. W rozpatrywanym zagadnieniu, Multimodalne obciążenie poznawcze i adaptacja interfejsu AR, kluczowe jest rozdzielenie właściwości systemu od zmienności zachowania użytkownika. Porównanie należy zaprojektować tak, aby mierzyć jednoznaczne zjawisko, a nie tylko ogólną ocenę technologii.

\section{Cel pracy i pytanie badawcze}

Czy prosta adaptacja zagęszczenia informacji HUD w AR ogranicza błędy względem stałego interfejsu? W minimalnym wariancie porównuje się jedną jawną regułę adaptacji z warunkiem statycznym, natomiast estymator wielomodalny stanowi odrębne rozszerzenie. Dla wyników o charakterze potwierdzającym konieczne jest zamrożenie protokołu.

\section{Zakres i odrębny wkład naukowy}
Przyczynowy wpływ prostej, jawnej reguły zagęszczania HUD na błędy wykonywania zadania. Klasyfikator obciążenia i fuzja biometryczna są odłożonym rozszerzeniem, nie częścią pilotażu P01. Dodatkowa wartość to efekt zdefiniowanej i powtarzalnej zmiany HUD, a nie ponowna deklaracja że AR zmniejsza obciążenie.

\section{Hipotezy badawcze}
\paragraph{H1.} Przy identycznej treści i trudności zadania adaptowany HUD prowadzi do niższego odsetka błędnych prób niż statyczny HUD.
\paragraph{H2.} Zmiana liczby informacji HUD zmniejsza liczbę niepożądanych przerwań poznawczych bez zwiększenia liczby pominiętych instrukcji.
\paragraph{H3.} Korzyść nie zanika po uwzględnieniu kolejności warunków i uprzedniego doświadczenia uczestnika z AR.

\section{Projekt badania i plan eksperymentów}
\subsection{A. Operacjonalizacja i kwalifikacja danych}
Przyczynowy wpływ prostej, jawnej reguły zagęszczania HUD na błędy wykonywania zadania. Klasyfikator obciążenia i fuzja biometryczna są odłożonym rozszerzeniem, nie częścią pilotażu P01. Jednostka wnioskowania: uczestnik × z góry ustalony scenariusz, z próbami zagnieżdżonymi w uczestniku. Punkt końcowy i sposób oceny będą zamrożone przed próbą potwierdzającą. Należy udokumentować pochodzenie obserwacji, czas akwizycji, identyfikatory sesji, harmonogram oraz legalność użycia materiałów.

\subsection{B. Eksperyment porównawczy}
Wewnątrzosobowy układ kontrbalansowany: ta sama liczba i jakość wskazówek, dwa poziomy HUD, wspólny czas ekspozycji. Estymand pierwszorzędowy: różnica marginalnego prawdopodobieństwa błędnej próby: adaptowany HUD minus statyczny HUD. Przypisanie warunków, mianowniki i kryteria wyłączeń nie mogą być dobierane po przejrzeniu wyników.

\subsection{C. Kontrole i analiza odporności}
Rozdzielić błąd zadania od samoopisu obciążenia; rejestrować skokowe zmiany HUD, liczbę informacji i rzeczywisty czas renderowania. Ocenić osoby o różnym doświadczeniu XR oraz proste i bardziej złożone kroki; nie włączać niezweryfikowanej okulometrii HMD. Każdą odmowę systemu, niekompletną próbę i brak sygnału należy raportować osobno, a nie po cichu usuwać z próby.

\subsection{D. Walidacja na niezależnym materiale}
Ocenę należy wykonać na uprzednio odłożonych scenach, sesjach lub uczestnikach. Hierarchiczny model logistyczny błędu próby z uczestnikiem i sceną jako czynnikami losowymi, przedziały dla różnicy ryzyka. Model i próg zatrzymać przed analizą końcową; rejestrować również brak przewagi i działania niepożądane.

\section{Estymand, analiza statystyczna i decyzje}
Główny estymand: różnica marginalnego prawdopodobieństwa błędnej próby: adaptowany HUD minus statyczny HUD. Jednostką powtarzalności jest uczestnik × z góry ustalony scenariusz, z próbami zagnieżdżonymi w uczestniku. Hierarchiczny model logistyczny błędu próby z uczestnikiem i sceną jako czynnikami losowymi, przedziały dla różnicy ryzyka. H1 jest kandydatem na hipotezę potwierdzającą. H2 i H3 pozostają eksploracyjne, dopóki protokół nie określi ich hierarchii i kontroli wielokrotnych porównań. Minimalny efekt istotny praktycznie, liczebność prób, margines nie gorszości i progi jakości wyznaczyć na podstawie pilotażu i analizy precyzji, nie zgadywać wartości.

\section{Stan badań i granica nowości}
\begin{itemize}
\item Existing adaptive physiological feedback in AR means adapting an interface is not by itself a novel contribution. \cite{xrprior1}.
\item Average reductions in learning load do not establish gains from the particular adaptive controller in a new scene. \cite{xrprior2}.
\end{itemize}
Dodatkowa wartość to efekt zdefiniowanej i powtarzalnej zmiany HUD, a nie ponowna deklaracja że AR zmniejsza obciążenie. Samo zastosowanie VLM, EEG, okulometrii lub adaptacyjnego HUD nie stanowi wystarczającej nowości międzynarodowej.

\section{Warunki bezpieczeństwa, sprzętu i zgodności}
Opis stanowi prospektywny plan, a nie prerejestrację ani zatwierdzoną procedurę. Przed zbieraniem danych od ludzi niezbędna jest właściwa weryfikacja etyczna i zgoda. Standardowy Quest 3 nie ma wbudowanego eye trackingu, a Tobii Eye Tracker 5 działa na stanowisku ekranowym. Odstępstwa, obciążenia i ograniczenia czujników muszą być ujawnione. Nie publikować surowych danych biometrycznych ani nagrań uczestników.

\section{Materiał i metody}

Przewidziano zrównoważone próby ze stałym HUD oraz z HUD adaptowanym według jednej jawnej reguły; rozbudowany model wielomodalny będzie osobnym rozszerzeniem. Wymagane jest zdefiniowanie jednostki analizy, zasad doboru przypadków, wersji scenariusza i źródeł pomiaru. Rozkład warunków oraz reguły wykluczenia prób należy określić przed pomiarem. W badaniach z udziałem ludzi niezbędna jest uprzednia właściwa ocena etyczna i świadoma zgoda. W testach wyłącznie symulacyjnych należy oddzielić prawdę symulacyjną od obserwacji empirycznej.

\subsection{Wynik główny i kryteria jakości}

Odsetek prób z błędnym wykonaniem zadania, przy jednoznacznym mianowniku. Pomiar powinien zachowywać jawne jednostki, znaczniki czasu i reguły oceny poprawności. Przy modelach predykcyjnych stosować rozłączne partycje oceny; przy eksperymentach interwencyjnych oddzielić skuteczność predykcji od przyczynowego wpływu interfejsu.

\section{Plan analizy}

Hierarchiczny model dwumianowy dla poprawności prób, kontrast adaptacyjny wobec stałego HUD; osobno badać trafność estymatora. W wynikach należy raportować wielkość efektu lub błąd estymacji z przedziałem niepewności, liczbą obserwacji i odsetkiem odrzuconych prób. Analizy dodatkowe wymagają jasnego oznaczenia jako eksploracyjne. Nie wybierać progów decyzyjnych ani modelu na podstawie obserwacji ze zbioru testowego.

\section{Wyniki}

W obecnej wersji nie raportuje się wyników pomiarów, tabel statystycznych ani parametrów wyznaczonych z nieudokumentowanych danych. Nie stwierdza się w tym opracowaniu, że eksperyment został wykonany lub zgłoszony do czasopisma.

\section{Dyskusja}

Interpretacja przyszłych wyników musi uwzględnić alternatywne mechanizmy i zależność wniosków od jakości danych, warunków zadania oraz ograniczeń sprzętu. Wartość potencjalnego rozwiązania zależy nie tylko od trafności modelu, lecz także od jego opóźnienia, stabilności i wpływu na zachowanie człowieka. Wyniki testów syntetycznych nie mogą być automatycznie uogólniane na populację uczestników.

\section{Ograniczenia badania}

Nie przypisywać przyczynowego efektu interfejsu samej trafności predykcji obciążenia; Quest 3 nie zapewnia okulometrii HMD. Granice nowości i trafności ekologicznej wymagają niezależnej weryfikacji. Żadna część szkicu nie dowodzi przeprowadzenia zbierania danych, uzyskania zgód etycznych ani gotowości manuskryptu do zgłoszenia.

\section{Wnioski}

Przedstawiono problem, kontrast i sposób jego planowanej oceny. Rozstrzygnięcie postawionego pytania wymaga przygotowania protokołu, zweryfikowanych źródeł oraz poprawnie udokumentowanych wyników. Bez tych materiałów nie można sformułować konkluzji empirycznych.

\begin{thebibliography}{9}
\bibitem{xrprior1} Designing adaptive feedback systems for managing cognitive load in augmented reality (2025), Design Science 11, e49. \url{https://doi.org/10.1017/dsj.2025.10040}.
\bibitem{xrprior2} Can augmented reality technology reduce learners' cognitive load? A meta-analysis (2026), Smart Learning Environments. \url{https://doi.org/10.1186/s40561-025-00429-7}.
\end{thebibliography}
\end{document}

